\documentclass[11pt]{article}
\usepackage{preamble}
\renewcommand{\answer}[1]{}

\begin{document}

\ladderheading{AP Statistics \textperiodcentered\ Unit 1 \textperiodcentered\ Topic 1.3 \textperiodcentered\ B: 2026-09-11 \textperiodcentered\ E: 2026-09-14}{A Percent Needs Its Base}

\focusbanner{TODAY'S PROBLEM}

Two clubs asked the student council for money to run a Saturday session. Each club surveyed \emph{all} of its members: ``Would you attend a Saturday session?'' The table shows every response as a count.\par\smallskip At the council meeting, the \textbf{Robotics club} presented a \textbf{frequency table} (counts) and the \textbf{Chess club} presented a \textbf{relative-frequency table} (percents). Each club chose the table that made its case look best.\par
\begin{center}
\textbf{``Would you attend a Saturday session?'' (every member surveyed; counts)}\par\smallskip
\begin{tabular}{|l|c|c|c|}
\hline
\textbf{Club} & \textbf{Yes} & \textbf{No} & \textbf{Total} \\ \hline
\textbf{Chess} & 15 & 3 & 18 \\ \hline
\textbf{Robotics} & 27 & 18 & 45 \\ \hline
\end{tabular}
\end{center}
\begin{firsttakebox}{FIRST TAKE --- before the video (5 min)}
Which club has more interest in a Saturday session --- Chess or Robotics? One sentence and one number.
\workspace{0.9in}
\end{firsttakebox}

\begin{callout}{\IconBook\ TABLES FOR A CATEGORICAL VARIABLE (keep on the page)}{calloutgreen}
A \textbf{frequency table} gives the count in each category. A \textbf{relative-frequency table} gives its proportion or percent: $\text{relative frequency}=\dfrac{\text{category count}}{\text{group total}}$. Multiply by 100 for a percent. Counts answer `how many?'; percents answer `what share?'. Always identify the group that a percent describes.
\end{callout}

\newpage
\textbf{Finish after the video:}\par\smallskip

\dokbadge{1}~\textbf{(a)} How many Robotics members said Yes?\par
\workspace{0.8in}

\dokbadge{2}~\textbf{(b)} Find the percent saying Yes in each club (nearest whole percent). Which is higher?\par
\workspace{1.4in}

\dokbadge{3}~\textbf{(c)} Which club has more student interest? State whether you mean more people or a larger share, and use the matching numbers. Explain why the other measure favors the other club and what a reader needs to know to compare them fairly.\par
\workspace{2.0in}

\begin{sentenceframebox}\raggedright\textbf{Frame:} By more interest, I mean \blank[1.5in], so I choose \blank[0.9in]: \blank[1.8in].\end{sentenceframebox}

\begin{sentenceframebox}\raggedright\textbf{Frame:} The other measure favors \blank[0.9in] because the club totals are \blank[1.4in]. Readers also need \blank[1.5in].\end{sentenceframebox}

\begin{summaryexitbox}{EXIT --- turn this in}
One thing the video changed about my first take: \hrulefill\par
\workspace{0.4in}
\end{summaryexitbox}

\end{document}
